\section{Prescribed-degree factors}\label{sec:f-factors}

Recall that an \emph{$f$-factor} of $G=([n],E)$ is a set $F\subseteq E$
with $\deg_F(v)=f(v)$ for every vertex $v$. We reduce its existence and
construction to one maximum-matching computation, adapting Tutte's
factor-to-perfect-matching reduction~\cite{Tutte1954Factor} with the split
network of Lemma~\ref{uniform:cardinality}. The underlying rearrangeable
network is the Bene\v{s} construction proved in
Lemma~\ref{uniform:network}~\cite{PermutationNetworks}.
In the incidence reduction, each vertex has one port per incident edge,
and $\deg_G(v)-f(v)$ of its ports must be matched locally. Complete
connections to that many auxiliary vertices can cost
$\deg_G(v)(\deg_G(v)-f(v))$ edges. A separate permutation network at
each vertex permits any such subset of ports at logarithmic expansion
in its degree.

\begin{corollary}[Almost-linear $f$-factor search]
\label{cor:f-factor}
There are a single uniform randomized word-RAM algorithm, fixed constants
$c_f\ge2$ and $C_f\ge1$, and a nonincreasing function
$\eta_f:\mathbb N_{\ge1}\to\mathbb R_{\ge0}$ with $\eta_f(q)\to0$ as
$q\to\infty$ such that the following holds.
The input is a simple graph $G=([n],E)$ with $n\ge1$, $m=|E|$, and
$p=n+m$, in the adjacency-list representation of Section~\ref{sec:introduction},
together with one-word integers $f(v)$ satisfying
$0\le f(v)\le\deg_G(v)$ for all $v$.
With the instruction set and costs of Section~\ref{sec:introduction} and
word length $w_f=\lceil c_f\log_2(n+2)\rceil$, the algorithm halts on every
computation path in at most
\[
                         C_f p^{1+\eta_f(p)}
\]
instructions. Every output is either an explicit edge list of a valid
$f$-factor or \textnormal{NO}. If an $f$-factor exists, it returns one with
probability at least $2/3$; otherwise it always returns \textnormal{NO}.
\end{corollary}

The $n$ degree records keep the input length within a constant factor of
$p$. The problem selects an edge subset; it does not encode binary edge
multiplicities or optimize edge weights.

\begin{lemma}[Linearithmic factor reduction]
\label{factor:reduction}
For the input of Corollary~\ref{cor:f-factor}, one can construct a possibly
empty simple graph $H_f$ with
$|V(H_f)|+|E(H_f)|=O(m\log(n+2))$
such that $H_f$ has a perfect matching if and only if $G$ has an $f$-factor.
The construction takes $O(n+m\log(n+2))$ word instructions and records.
A perfect matching of $H_f$ yields an explicit $f$-factor within the same
bound.
\end{lemma}

\begin{proof}
For every host edge $e=\{u,v\}$, create distinct port vertices
$x_{u,e}$ and $x_{v,e}$ and join them by a \emph{cross edge} labeled by $e$.
For a vertex $v$ write $d_v=\deg_G(v)$ and $b_v=d_v-f(v)$. If $d_v=0$,
then $f(v)=0$ and no local construction is needed. Otherwise, let $s_v$ be
the least power of two at least $d_v$ and take a separate copy of the
recursive network in Lemma~\ref{uniform:network}. Assign distinct input
positions to the $d_v$ incidences at $v$, and select $b_v$ distinct output
positions.

Split every network position $a$ into vertices $a_{\mathrm{in}}$ and
$a_{\mathrm{out}}$ joined by an idle edge. For every directed arc $a\to a'$
of the network, add the transition edge
$a_{\mathrm{out}}a'_{\mathrm{in}}$. Join $x_{v,e}$ to the input endpoint
of its assigned input position. At each selected output position, add a new
sink vertex adjacent only to that position's output endpoint. All these
local constructions use disjoint new vertices. The recursive network has
distinct arcs, and the edge types have distinct endpoint pairs, so $H_f$ is
simple. This also covers $s_v=1$, when the single position is both an input
and an output and there is no transition edge.

Consider any perfect matching of $H_f$. At every split position, either its
idle edge is used, or both endpoints are matched by non-idle edges. Therefore
among endpoints matched by non-idle edges, the input and output counts in
the network at $v$ are equal. Each selected transition edge contributes once
to each number and
cancels from this equality. A port attachment contributes only an input
endpoint, and a sink attachment contributes only an output endpoint. All
$b_v$ degree-one sinks must be covered, so exactly $b_v$ port attachments
are used. Every port has only its attachment and its cross edge. Hence
exactly $d_v-b_v=f(v)$ ports at $v$ use cross edges. The labels of all
selected cross edges form a set $F\subseteq E$ with $\deg_F(v)=f(v)$ for
every $v$.

Conversely, suppose that $F$ is an $f$-factor. At $v$, exactly $b_v$ ports
belong to edges outside $F$. Map their input positions bijectively to the
selected outputs and extend this map to a permutation of all $s_v$ input
and output positions. Lemma~\ref{uniform:network} routes this permutation
by vertex-disjoint directed paths; retain only the $b_v$ paths just chosen.
Select the cross edge for each $e\in F$. Along each retained path, select
its port attachment, all its transition edges, and its final sink edge.
Each split position on these paths is covered once at its input endpoint
and once at its output endpoint. The paths are vertex-disjoint, so these
choices do not conflict. Use idle edges at every unused position. All ports,
split endpoints, and sinks are now covered exactly once. When $b_v=0$ no
path is retained; a length-zero path when $s_v=1$ uses just its port and
sink attachments. Thus the construction gives a perfect matching of $H_f$.

For $d_v>0$ we have $s_v\le2d_v$, so the local network has
$O(d_v\log(d_v+2))$ positions and arcs. Ports, cross edges, attachments,
and sinks add $O(m)$ more records. Since $d_v\le n-1$ and
$\sum_v d_v=2m$, their total is $O(m\log(n+2))$. The uniform network
construction and a scan of the adjacency lists take
$O(n+m\log(n+2))$ instructions. Scanning the occurrence with the smaller
endpoint label assigns an identifier to each host edge without assuming any
list order. Two passes over the constructed edge list count degrees and fill
the adjacency-list representation of $H_f$ within the same bound.
Every identifier uses $O(\log(n+2))$
bits because the constructed size is polynomial in $n$. Store each cross
edge's host identifier; scanning a returned perfect matching and selecting
these labels gives the stated decoder bound.
\end{proof}

\begin{proof}[Proof of Corollary~\ref{cor:f-factor}]
If $m=0$, validity forces every $f(v)=0$, so return the empty edge list.
Otherwise $H_f$ is nonempty; construct it and apply the single algorithm of
Theorem~\ref{thm:main}. Check the returned edge identifiers and one-time
coverage of every vertex, returning \textnormal{NO} unless the list is a
perfect matching. Section~\ref{sec:uniformization}'s output implementation
copies at most $|V(H_f)|/2$ edges on every path, so the check takes
$O(|V(H_f)|+|E(H_f)|)$ instructions, including initialization. If it passes,
decode the cross-edge labels by Lemma~\ref{factor:reduction}, checking their
host identifiers and prescribed degrees.

When the matching call returns a maximum matching, the test accepts exactly
when $H_f$ has a perfect matching and then supplies the factor by the
reduction. This maximum-matching event has probability at least $2/3$.
The explicit checks ensure that every accepted edge list is a valid factor,
regardless of the matching call's outcome. Thus a failed call can cause a
false negative, whereas an input with no factor always yields \textnormal{NO}.

Put $P=|V(H_f)|+|E(H_f)|=O(p\log(p+2))$.
Since $|V(H_f)|$ is polynomial in $n$, the matching algorithm's virtual
word length $\lceil c_0\log_2(|V(H_f)|+2)\rceil$ is at most
$\lceil c_f\log_2(n+2)\rceil$ for one fixed $c_f$. Choose $c_f$ large
enough also for the construction's indices. Simulate each virtual word
instruction by a constant number of physical instructions, masking to the
virtual width for modular arithmetic and Boolean operations. Quotients,
remainders, comparisons, shifts, and memory addresses then have the same
semantics as in the matching call; each fair random bit remains charged.
The fixed masks and input conversion take at most the reduction's work
bound. Thus this is one finite program with a fixed physical word constant,
not a separate algorithm for each input size or degree prescription.

For fixed $\xi>0$, Theorem~\ref{thm:main} gives $O_\xi(P^{1+\xi/2})$
every-path instructions, absorbing small $P$ into the constant. Together
with all reduction, simulation, checking, and output work, the bound on $P$
gives $O_\xi(p^{1+\xi})$ for large $p$, for this same program for
every $\xi>0$. Boundedness of $\eta$ also gives a fixed polynomial every-path
bound. The supremum-envelope construction at the end of
Section~\ref{sec:uniformization} now gives a nonnegative nonincreasing
$\eta_f(q)\to0$; a fixed $C_f$ covers the small input sizes, including $p=1$.
\end{proof}
